\documentclass[11pt,a4paper]{article}

% ── packages ──────────────────────────────────────────────────────────────
\usepackage[utf8]{inputenc}
\usepackage[T1]{fontenc}
\usepackage{amsmath,amssymb,bm}
\usepackage{geometry}
\geometry{margin=2.5cm}
\usepackage{booktabs}
\usepackage{hyperref}
\usepackage{algorithm}
\usepackage{algpseudocode}
\usepackage{graphicx}
\usepackage{xcolor}
\usepackage{mathtools}
\usepackage{cleveref}

% ── macros ────────────────────────────────────────────────────────────────
\newcommand{\se}{\mathfrak{se}(3)}
\newcommand{\SO}{\mathrm{SO}(3)}
\newcommand{\SE}{\mathrm{SE}(3)}
\newcommand{\Ad}{\mathrm{Ad}}
\newcommand{\ad}{\mathrm{ad}}
\newcommand{\Jac}{\mathcal{J}}
\newcommand{\curly}[1]{\mathcal{S}\!\left(#1\right)}
\newcommand{\R}{\mathbb{R}}
\newcommand{\inv}{^{-1}}
\newcommand{\T}{\mathbf{T}}
\newcommand{\bxi}{\bm{\xi}}
\newcommand{\bvarpi}{\bm{\varpi}}
\newcommand{\bgamma}{\bm{\gamma}}
\newcommand{\bPhi}{\bm{\Phi}}
\newcommand{\be}{\mathbf{e}}
\newcommand{\bF}{\mathbf{F}}
\newcommand{\bA}{\mathbf{A}}
\newcommand{\bB}{\mathbf{B}}
\newcommand{\bC}{\mathbf{C}}
\newcommand{\bD}{\mathbf{D}}
\newcommand{\bQ}{\mathbf{Q}}
\newcommand{\bL}{\mathbf{L}}
\newcommand{\bI}{\mathbf{I}}
\newcommand{\bzero}{\mathbf{0}}
\DeclareMathOperator{\diag}{diag}

\title{Extending the Multi-Robot Continuum State Estimator\\
with General Transition Functions and Control Inputs}
\author{}
\date{}

\begin{document}
\maketitle

\begin{abstract}
This document provides a detailed mathematical description of the extensions made to the multi-robot continuum state estimator of~\cite{lilge2022multi} (arXiv: 2210.14842) to incorporate general transition functions and control inputs, following the formulation of~\cite{lilge2024pcr} (arXiv: 2408.01333).
The original estimator hardcodes the White-Noise-on-Acceleration (WNOA) Gaussian process prior, which assumes a linear, time-invariant (LTI) transition function and analytically known covariance.
The extension replaces these closed-form WNOA expressions with configurable, input-dependent transition functions $\bPhi$, transition integrals $\bPhi_{\mathrm{int}}$, and covariance matrices $\bQ$ that are computed via the Magnus expansion on $\SE$.
This generalisation affects four core components of the estimator: the prior factor assembly, the prior cost evaluation, the GP interpolation scheme, and the initial-guess construction.
\end{abstract}

%==========================================================================
\section{Background: State Representation and Estimation Problem}
\label{sec:background}
%==========================================================================

\subsection{State at Each Node}
Consider $N$ continuum robots, each discretised into $K_n$ estimation nodes along arclength~$s$.
At each node~$k$ of robot~$n$, the state consists of:
\begin{itemize}
    \item A homogeneous transformation $\T_{b_k i} \in \SE$ relating the body frame at node~$k$ to the inertial frame.
    \item A generalised strain vector $\bvarpi_k \in \R^6$ combining translational strains $\bm{\nu}$ and rotational strains (curvatures) $\bm{\omega}$.
\end{itemize}

The optimisation operates internally in the body-to-inertial convention $\T_{b_k i}$.
The local state increment between nodes $k$ and $k{+}1$ is captured by
\begin{equation}
    \bxi_{k,k+1} = \ln\!\bigl(\T_{b_{k+1} i}\,\T_{b_k i}\inv\bigr)^\vee \in \R^6,
\end{equation}
and the compound 12-dimensional local state at node~$k$ is defined as
\begin{equation}
    \bgamma_k = \begin{pmatrix} \bzero_{6\times 1} \\ \bvarpi_k \end{pmatrix}, \qquad
    \bgamma_{k+1} = \begin{pmatrix} \bxi_{k,k+1} \\ \Jac\inv(\bxi_{k,k+1})\,\bvarpi_{k+1} \end{pmatrix},
\end{equation}
where $\Jac\inv(\bxi)$ is the inverse of the left Jacobian of $\SE$ evaluated at $\bxi$.

\subsection{MAP Estimation Problem}
The estimator solves a maximum a posteriori (MAP) problem by iteratively linearising a nonlinear least-squares objective:
\begin{equation}
    \min_{\mathbf{x}} \;\; \underbrace{\sum_{\text{prior}} \tfrac{1}{2}\,\be_p^\top \bQ\inv \be_p}_{\text{smoothness}} \;+\; \underbrace{\sum_{\text{coupling}} \tfrac{1}{2}\,\be_c^\top \mathbf{R}_c\inv \be_c}_{\text{kinematic constraints}} \;+\; \underbrace{\sum_{\text{meas.}} \tfrac{1}{2}\,\be_m^\top \mathbf{R}_m\inv \be_m}_{\text{observations}},
    \label{eq:map}
\end{equation}
where $\mathbf{x}$ collects all node states and (optionally) a common end-effector pose.
Each iteration assembles the Gauss-Newton normal equations
$\mathbf{H}\,\delta\mathbf{x} = \mathbf{g}$
and applies a state update (optionally with backtracking line search).

%==========================================================================
\section{Original Formulation: Hardcoded WNOA Prior}
\label{sec:original}
%==========================================================================

The original estimator~\cite{lilge2022multi} uses the White-Noise-on-Acceleration (WNOA) Gaussian process prior on $\SE$.
This models the arclength evolution of the robot backbone as a second-order stochastic differential equation driven by white noise in the strain-rate space.

\subsection{WNOA Transition Function}
\label{sec:wnoa-phi}
The WNOA transition function over an arclength interval $\Delta s = s_{k+1} - s_k$ is the linear, time-invariant (LTI) matrix:
\begin{equation}
    \bPhi_{\mathrm{WNOA}}(\Delta s) = \begin{pmatrix} \bI_6 & \Delta s\,\bI_6 \\ \bzero_6 & \bI_6 \end{pmatrix} \in \R^{12 \times 12}.
    \label{eq:phi-wnoa}
\end{equation}

\subsection{WNOA Inverse Covariance}
\label{sec:wnoa-qinv}
The corresponding inverse covariance (information) matrix is the analytically known block matrix:
\begin{equation}
    \bQ_{\mathrm{WNOA}}\inv(\Delta s) = \begin{pmatrix}
        \dfrac{12}{\Delta s^3}\,\bQ_c\inv & -\dfrac{6}{\Delta s^2}\,\bQ_c\inv \\[8pt]
        -\dfrac{6}{\Delta s^2}\,\bQ_c\inv & \dfrac{4}{\Delta s}\,\bQ_c\inv
    \end{pmatrix} \in \R^{12 \times 12},
    \label{eq:qinv-wnoa}
\end{equation}
where $\bQ_c \in \R^{6 \times 6}$ is the continuous-time power spectral density of the strain-driving noise.
There is no transition integral term: $\bPhi_{\mathrm{int}} = \bzero_{12 \times 1}$.

\subsection{WNOA Prior Residual and Jacobian}
\label{sec:wnoa-prior}
The prior residual between nodes $k$ and $k{+}1$ is:
\begin{equation}
    \be_{\mathrm{WNOA}} = \begin{pmatrix} \bxi_{k,k+1} - \Delta s\,\bvarpi_k \\ \Jac\inv(\bxi_{k,k+1})\,\bvarpi_{k+1} - \bvarpi_k \end{pmatrix} \in \R^{12}.
    \label{eq:e-wnoa}
\end{equation}
Note that this is exactly $\bgamma_{k+1} - \bPhi_{\mathrm{WNOA}}\,\bgamma_k$.

The Jacobian $\bF \in \R^{12 \times 24}$ with respect to the stacked state $(\bgamma_k, \bgamma_{k+1})$ is:
\begin{equation}
    \bF_{\mathrm{WNOA}} = \begin{pmatrix}
        -\Jac\inv \Ad & -\Delta s\,\bI_6 & \Jac\inv & \bzero_6 \\
        -\tfrac{1}{2}\curly{\bvarpi_{k+1}}\Jac\inv \Ad & -\bI_6 & \tfrac{1}{2}\curly{\bvarpi_{k+1}}\Jac\inv & \Jac\inv
    \end{pmatrix},
    \label{eq:F-wnoa}
\end{equation}
where $\Ad = \Ad(\T_{b_{k+1}i}\,\T_{b_k i}\inv)$, $\Jac\inv = \Jac\inv(\bxi_{k,k+1})$, and $\curly{\cdot}$ denotes the curly-hat operator.

\subsection{WNOA GP Interpolation}
\label{sec:wnoa-interp}
Interpolation between estimation nodes uses the GP conditional mean.
For a query arclength $s_q \in (s_k, s_{k+1})$ with $\Delta s_q = s_q - s_k$, the WNOA interpolation matrices are:
\begin{align}
    \bQ(s_q) &= \begin{pmatrix} \frac{\Delta s_q^3}{3}\bQ_c & \frac{\Delta s_q^2}{2}\bQ_c \\ \frac{\Delta s_q^2}{2}\bQ_c & \Delta s_q\,\bQ_c \end{pmatrix}, \label{eq:Q-interp-wnoa}\\[4pt]
    \bPhi(s_q, s_{k+1}) &= \begin{pmatrix} \bI_6 & (s_{k+1}-s_q)\bI_6 \\ \bzero_6 & \bI_6 \end{pmatrix}, \label{eq:phi-partial-wnoa}\\[4pt]
    \bPhi(s_k, s_q) &= \begin{pmatrix} \bI_6 & \Delta s_q\,\bI_6 \\ \bzero_6 & \bI_6 \end{pmatrix}. \label{eq:lambda-wnoa}
\end{align}
The GP interpolation coefficients are then:
\begin{equation}
    \bm{\Psi} = \bQ(s_q)\,\bPhi(s_q, s_{k+1})^\top\,\bQ_{\mathrm{WNOA}}\inv, \qquad
    \bm{\Lambda} = \bPhi(s_k, s_q) - \bm{\Psi}\,\bPhi_{\mathrm{WNOA}},
\end{equation}
and the interpolated state is:
\begin{equation}
    \bgamma_q = \bm{\Lambda}\,\bgamma_k + \bm{\Psi}\,\bgamma_{k+1}.
    \label{eq:interp-wnoa}
\end{equation}

%==========================================================================
\section{Extended Formulation: General Transition Functions}
\label{sec:extended}
%==========================================================================

The extension replaces the hardcoded WNOA expressions with configurable transition functions parameterised by a \emph{control input} $\mathbf{u}(s)$ on each segment.
The control input is a 12-dimensional vector
\begin{equation}
    \mathbf{u}(s) = \begin{pmatrix} \bm{\nu}_u(s) \\ \bm{\omega}_u(s) \\ \dot{\bm{\nu}}_u(s) \\ \dot{\bm{\omega}}_u(s) \end{pmatrix} \in \R^{12},
\end{equation}
combining velocity (strain) and acceleration components in $\se$.
Three control input types are supported:

\begin{description}
    \item[\texttt{None}] No control input. The standard WNOA prior is recovered.
    \item[\texttt{Constant}] The control input $\mathbf{u}_0$ is constant over the segment $[s_k, s_{k+1}]$.
    \item[\texttt{PiecewiseLinear}] The control input varies linearly between $P$ waypoints $\mathbf{u}_0, \mathbf{u}_1, \ldots, \mathbf{u}_{P-1}$, defining $P{-}1$ subsegments within each estimation interval.
\end{description}

%--------------------------------------------------------------------------
\subsection{General Transition Function \texorpdfstring{$\bPhi$}{Phi}}
\label{sec:general-phi}
%--------------------------------------------------------------------------

\subsubsection{Type: None (WNOA Recovery)}
When no control input is present, the transition function reduces to the WNOA form~\eqref{eq:phi-wnoa}:
\begin{equation}
    \bPhi_{\texttt{None}}(\Delta s) = \begin{pmatrix} \bI_6 & \Delta s\,\bI_6 \\ \bzero_6 & \bI_6 \end{pmatrix}.
\end{equation}

\subsubsection{Type: Constant}
\label{sec:phi-constant}
Given a constant control input $\mathbf{u}_0 \in \R^{12}$, define the \emph{system matrix}
\begin{equation}
    \bA = \begin{pmatrix} \tfrac{1}{2}\curly{\mathbf{v}} & \bI_6 \\ \tfrac{1}{2}\curly{\mathbf{a}} & -\tfrac{1}{2}\curly{\mathbf{v}} \end{pmatrix} \in \R^{12 \times 12},
    \label{eq:A-constant}
\end{equation}
where $\mathbf{v} = -\mathbf{u}_0^{[1:6]}$ (velocity) and $\mathbf{a} = -\mathbf{u}_0^{[7:12]}$ (acceleration), and $\curly{\cdot}$ is the $6{\times}6$ curly-hat operator on $\se$.

The transition function is approximated via 5th-order matrix Taylor expansion (truncated matrix exponential):
\begin{equation}
    \bPhi_{\texttt{Const}}(\Delta s) = \sum_{j=0}^{5} \frac{\bA^j\,\Delta s^j}{j!} = \bI + \bA\,\Delta s + \frac{\bA^2}{2}\Delta s^2 + \frac{\bA^3}{6}\Delta s^3 + \frac{\bA^4}{24}\Delta s^4 + \frac{\bA^5}{120}\Delta s^5.
    \label{eq:phi-constant}
\end{equation}

\subsubsection{Type: PiecewiseLinear --- Magnus Expansion}
\label{sec:phi-pwl}

For piecewise-linear inputs, each subsegment $i$ has endpoints $\mathbf{u}_i$ and $\mathbf{u}_{i+1}$.
Three sub-cases arise:

\paragraph{Sub-case 1: Both endpoints near zero.}
If $\|\mathbf{u}_i\| < \epsilon$ and $\|\mathbf{u}_{i+1}\| < \epsilon$, the WNOA transition is used for that subsegment.

\paragraph{Sub-case 2: Constant input over subsegment.}
If $\mathbf{u}_i \approx \mathbf{u}_{i+1}$, the constant-input formula~\eqref{eq:phi-constant} is used with $\tau = \Delta s_{\mathrm{sub}}$.

\paragraph{Sub-case 3: General linear variation.}
Define the linear decomposition $\mathbf{u}(s) = \mathbf{u}_{(0)} + s\,\mathbf{u}_{(1)}$ where
$\mathbf{u}_{(0)} = -\mathbf{u}_i$, \;
$\mathbf{u}_{(1)} = -(\mathbf{u}_{i+1} - \mathbf{u}_i)/\Delta s_{\mathrm{sub}}$,
and construct:
\begin{align}
    \bA_0 &= \begin{pmatrix} \tfrac{1}{2}\curly{\mathbf{u}_{(0)}^{[1:6]}} & \bI_6 \\ \tfrac{1}{2}\curly{\mathbf{u}_{(0)}^{[7:12]}} & -\tfrac{1}{2}\curly{\mathbf{u}_{(0)}^{[1:6]}} \end{pmatrix}, \label{eq:A0}\\[4pt]
    \bA_1 &= \begin{pmatrix} \tfrac{1}{2}\curly{\mathbf{u}_{(1)}^{[1:6]}} & \bzero_6 \\ \tfrac{1}{2}\curly{\mathbf{u}_{(1)}^{[7:12]}} & -\tfrac{1}{2}\curly{\mathbf{u}_{(1)}^{[1:6]}} \end{pmatrix}. \label{eq:A1}
\end{align}

Using the \textbf{Magnus expansion} with Lie brackets $[\cdot,\cdot]$, define the coefficients:
\begin{equation}
    \bA = \bA_0, \quad \bB = \tfrac{1}{2}\bA_1, \quad \bC = \tfrac{1}{12}[\bA_1, \bA_0], \quad \bD = \tfrac{1}{240}[\bA_1, [\bA_1, \bA_0]],
    \label{eq:magnus-coeff}
\end{equation}
where $[\bA_1, \bA_0] = \bA_1 \bA_0 - \bA_0 \bA_1$.

The 5th-order Magnus expansion for the subsegment transition function is:
\begin{align}
    \bPhi_{\mathrm{sub}} &= \bI + \bA\tau + \bigl(\tfrac{1}{2}\bA^2 + \bB\bigr)\tau^2 \nonumber\\
    &\quad + \bigl(\bC + \tfrac{1}{2}\bA\bB + \tfrac{1}{2}\bB\bA + \tfrac{1}{6}\bA^3\bigr)\tau^3 \nonumber\\
    &\quad + \bigl(\tfrac{1}{2}\bA\bC + \tfrac{1}{2}\bB^2 + \tfrac{1}{2}\bC\bA + \tfrac{1}{6}\bA^2\bB + \tfrac{1}{6}\bA\bB\bA + \tfrac{1}{6}\bB\bA^2 + \tfrac{1}{24}\bA^4\bigr)\tau^4 \nonumber\\
    &\quad + \bigl(\bD + \tfrac{1}{2}\bB\bC + \tfrac{1}{2}\bC\bB + \tfrac{1}{6}\bA^2\bC + \tfrac{1}{6}\bA\bB^2 + \tfrac{1}{6}\bA\bC\bA + \tfrac{1}{6}\bB\bA\bB + \tfrac{1}{6}\bB^2\bA \nonumber\\
    &\qquad\quad + \tfrac{1}{6}\bC\bA^2 + \tfrac{1}{24}\bA^3\bB + \tfrac{1}{24}\bA^2\bB\bA + \tfrac{1}{24}\bA\bB\bA^2 + \tfrac{1}{24}\bB\bA^3 + \tfrac{1}{120}\bA^5\bigr)\tau^5, \label{eq:phi-magnus}
\end{align}
where $\tau = \Delta s_{\mathrm{sub}}$.

The total transition function over the full segment is the ordered product of subsegment transitions:
\begin{equation}
    \bPhi_{\texttt{PwL}} = \bPhi_{\mathrm{sub}}^{(P-2)} \cdots \bPhi_{\mathrm{sub}}^{(1)}\,\bPhi_{\mathrm{sub}}^{(0)}.
\end{equation}

%--------------------------------------------------------------------------
\subsection{Transition Function Integral \texorpdfstring{$\bPhi_{\mathrm{int}}$}{Phi\_int}}
\label{sec:general-phi-int}
%--------------------------------------------------------------------------

The transition function integral captures the deterministic drift due to the control input.

\subsubsection{Type: None}
\begin{equation}
    \bPhi_{\mathrm{int, None}} = \bzero_{12 \times 1}.
\end{equation}

\subsubsection{Type: Constant}
Define the Taylor-expansion coefficients of $\bPhi$:
\begin{equation}
    \bPhi^{(j)} = \frac{\bA^j}{j!}, \quad j=1,\ldots,5.
    \label{eq:phi-taylor-coeff}
\end{equation}
The integral is:
\begin{equation}
    \bPhi_{\mathrm{int, Const}} = \mathbf{u}_0'\,\tau + \tfrac{1}{2}\bPhi^{(1)}\mathbf{u}_0'\,\tau^2 + \tfrac{1}{3}\bPhi^{(2)}\mathbf{u}_0'\,\tau^3 + \tfrac{1}{4}\bPhi^{(3)}\mathbf{u}_0'\,\tau^4 + \tfrac{1}{5}\bPhi^{(4)}\mathbf{u}_0'\,\tau^5 + \tfrac{1}{6}\bPhi^{(5)}\mathbf{u}_0'\,\tau^6,
    \label{eq:phi-int-constant}
\end{equation}
where $\mathbf{u}_0' = -\mathbf{u}_0$ and $\tau = \Delta s$.

\subsubsection{Type: PiecewiseLinear}
For each subsegment with linearly varying input, define the \emph{Magnus-corrected} expansion coefficients (note: these differ from the constant-input case because the Magnus correction terms $\bC$ and $\bD$ contribute at higher orders):
\begin{align}
    \bPhi^{(1)} &= \bA, \nonumber\\
    \bPhi^{(2)} &= \tfrac{1}{2}\bA^2, \nonumber\\
    \bPhi^{(3)} &= \bC + \tfrac{1}{6}\bA^3, \nonumber\\
    \bPhi^{(4)} &= \tfrac{1}{2}\bA\bC + \tfrac{1}{2}\bC\bA + \tfrac{1}{24}\bA^4, \nonumber\\
    \bPhi^{(5)} &= \bD + \tfrac{1}{6}\bA^2\bC + \tfrac{1}{6}\bA\bC\bA + \tfrac{1}{6}\bC\bA^2 + \tfrac{1}{120}\bA^5,
    \label{eq:phi-magnus-coeff}
\end{align}
where $\bC = \frac{1}{12}[\bA_1, \bA_0]$ and $\bD = \frac{1}{240}[\bA_1, [\bA_1, \bA_0]]$.
These are the order-$j$ coefficients of the $\bA$-dependent (i.e., $\bB$-independent) part of the Magnus expansion~\eqref{eq:phi-magnus}.

Additionally, define the cross-term expansion coefficients incorporating $\bB$:
\begin{align}
    \bPhi_{11} &= \bB, \nonumber\\
    \bPhi_{21} &= \tfrac{1}{2}(\bA\bB + \bB\bA), \quad \bPhi_{22} = \tfrac{1}{2}\bB^2, \nonumber\\
    \bPhi_{31} &= \tfrac{1}{6}(\bA^2\bB + \bA\bB\bA + \bB\bA^2), \quad \bPhi_{32} = \tfrac{1}{6}(\bA\bB^2 + \bB\bA\bB + \bB^2\bA), \nonumber\\
    \bPhi_{41} &= \tfrac{1}{2}(\bB\bC + \bC\bB) + \tfrac{1}{24}(\bA^3\bB + \bA^2\bB\bA + \bA\bB\bA^2 + \bB\bA^3).
    \label{eq:phi-mixed}
\end{align}

The subsegment integral (to 7th order in $\tau$) is:
\begin{align}
    \bPhi_{\mathrm{int,sub}} &= \mathbf{u}_{(0)}\tau + \tfrac{1}{2}\bigl(\bPhi^{(1)}\mathbf{u}_{(0)} + \mathbf{u}_{(1)}\bigr)\tau^2 \nonumber\\
    &\quad + \tfrac{1}{3}\bigl(\tfrac{1}{2}\bPhi^{(1)}\mathbf{u}_{(1)} + \bPhi^{(2)}\mathbf{u}_{(0)} + 2\bPhi_{11}\mathbf{u}_{(0)}\bigr)\tau^3 \nonumber\\
    &\quad + \tfrac{1}{4}\bigl(\tfrac{1}{3}\bPhi^{(2)}\mathbf{u}_{(1)} + \bPhi^{(3)}\mathbf{u}_{(0)} + \bPhi_{11}\mathbf{u}_{(1)} + \tfrac{5}{3}\bPhi_{21}\mathbf{u}_{(0)}\bigr)\tau^4 \nonumber\\
    &\quad + \tfrac{1}{5}\bigl(\tfrac{1}{4}\bPhi^{(3)}\mathbf{u}_{(1)} + \bPhi^{(4)}\mathbf{u}_{(0)} + \tfrac{7}{12}\bPhi_{21}\mathbf{u}_{(1)} + \tfrac{8}{3}\bPhi_{22}\mathbf{u}_{(0)} + \tfrac{3}{2}\bPhi_{31}\mathbf{u}_{(0)}\bigr)\tau^5 \nonumber\\
    &\quad + \tfrac{1}{6}\bigl(\tfrac{1}{5}\bPhi^{(4)}\mathbf{u}_{(1)} + \bPhi^{(5)}\mathbf{u}_{(0)} + \bPhi_{22}\mathbf{u}_{(1)} + \tfrac{2}{5}\bPhi_{31}\mathbf{u}_{(1)} + \tfrac{11}{5}\bPhi_{32}\mathbf{u}_{(0)} + \tfrac{7}{5}\bPhi_{41}\mathbf{u}_{(0)}\bigr)\tau^6 \nonumber\\
    &\quad + \tfrac{1}{7}\bigl(\tfrac{1}{6}\bPhi^{(5)}\mathbf{u}_{(1)} + \tfrac{19}{30}\bPhi_{32}\mathbf{u}_{(1)} + \tfrac{3}{10}\bPhi_{41}\mathbf{u}_{(1)}\bigr)\tau^7.
    \label{eq:phi-int-pwl}
\end{align}

The total integral is assembled by propagating each subsegment's contribution through subsequent transition functions:
\begin{equation}
    \bPhi_{\mathrm{int}} = \sum_{i=0}^{P-2} \left(\prod_{j=i+1}^{P-2} \bPhi_{\mathrm{sub}}^{(j)}\right) \bPhi_{\mathrm{int,sub}}^{(i)}.
\end{equation}

%--------------------------------------------------------------------------
\subsection{General Covariance \texorpdfstring{$\bQ$}{Q} and Information Matrix \texorpdfstring{$\bQ\inv$}{Q inv}}
\label{sec:general-qinv}
%--------------------------------------------------------------------------

Define the projection matrix
\begin{equation}
    \bL = \begin{pmatrix} \bzero_{6\times 6} \\ \bI_6 \end{pmatrix} \in \R^{12 \times 6}, \qquad
    \bL\bQ_c\bL^\top = \begin{pmatrix} \bzero_6 & \bzero_6 \\ \bzero_6 & \bQ_c \end{pmatrix} \in \R^{12 \times 12}.
    \label{eq:LQLT}
\end{equation}

\subsubsection{Type: None}
For WNOA, the covariance has the closed-form inverse~\eqref{eq:qinv-wnoa}, equivalent to:
\begin{equation}
    \bQ_{\texttt{None}} = \begin{pmatrix} \frac{\Delta s^3}{3}\bQ_c & \frac{\Delta s^2}{2}\bQ_c \\ \frac{\Delta s^2}{2}\bQ_c & \Delta s\,\bQ_c \end{pmatrix}.
\end{equation}

\subsubsection{Type: Constant}
\label{sec:Q-constant}
Using the expansion coefficients $\bPhi^{(j)} = \bA^j/j!$, the covariance is computed via a Taylor series of the integral $\int_0^\tau \bPhi(s)\,\bL\bQ_c\bL^\top\,\bPhi(s)^\top\,ds$:
\begin{align}
    \bQ_{\texttt{Const}} &= \bL\bQ_c\bL^\top\,\tau \nonumber\\
    &\quad + \tfrac{1}{2}\bigl(\bL\bQ_c\bL^\top{\bPhi^{(1)}}^\top + \bPhi^{(1)}\bL\bQ_c\bL^\top\bigr)\tau^2 \nonumber\\
    &\quad + \tfrac{1}{3}\bigl(\bL\bQ_c\bL^\top{\bPhi^{(2)}}^\top + \bPhi^{(2)}\bL\bQ_c\bL^\top + \bPhi^{(1)}\bL\bQ_c\bL^\top{\bPhi^{(1)}}^\top\bigr)\tau^3 \nonumber\\
    &\quad + \tfrac{1}{4}\bigl(\bL\bQ_c\bL^\top{\bPhi^{(3)}}^\top + \bPhi^{(3)}\bL\bQ_c\bL^\top + \bPhi^{(1)}\bL\bQ_c\bL^\top{\bPhi^{(2)}}^\top + \bPhi^{(2)}\bL\bQ_c\bL^\top{\bPhi^{(1)}}^\top\bigr)\tau^4 \nonumber\\
    &\quad + \tfrac{1}{5}\bigl(\ldots\text{symmetric pattern continues with } \bPhi^{(4)}\ldots\bigr)\tau^5 \nonumber\\
    &\quad + \tfrac{1}{6}\bigl(\ldots\text{symmetric pattern continues with } \bPhi^{(5)}\ldots\bigr)\tau^6.
    \label{eq:Q-constant}
\end{align}
The general pattern at order $p$ is:
\begin{equation}
    \bQ^{(p)} = \frac{1}{p}\sum_{\substack{i+j = p-1 \\ i,j \geq 0}} \bPhi^{(i)}\,\bL\bQ_c\bL^\top\,{\bPhi^{(j)}}^\top.
\end{equation}
The information matrix is then $\bQ\inv = \bQ^{-1}$, computed via dense matrix inversion.

\subsubsection{Type: PiecewiseLinear}
\label{sec:Q-pwl}

For PiecewiseLinear inputs, each subsegment contributes a covariance block $\bQ_{\mathrm{sub}}^{(i)}$.
The base pattern follows the same structure as the Constant case~\eqref{eq:Q-constant}, but using the Magnus-corrected coefficients $\bPhi^{(j)}$ from~\eqref{eq:phi-magnus-coeff} instead of $\bA^j/j!$.

When the input is linearly varying within a subsegment, additional cross-terms with $\bPhi_{11}, \bPhi_{21}, \bPhi_{22}, \bPhi_{31}, \bPhi_{32}, \bPhi_{41}$ appear.
The coefficients for the cross-terms are (listing only the additional terms beyond the $\bPhi^{(j)}$-only pattern):
\begin{align}
    \text{Order 3:}&\quad +\;\tfrac{2}{3}\bigl(\bL\bQ_c\bL^\top\bPhi_{11}^\top + \bPhi_{11}\bL\bQ_c\bL^\top\bigr)\tau^3 \nonumber\\
    \text{Order 4:}&\quad +\;\tfrac{5}{12}\bigl(\bL\bQ_c\bL^\top\bPhi_{21}^\top + \bPhi_{21}\bL\bQ_c\bL^\top + \bPhi^{(1)}\bL\bQ_c\bL^\top\bPhi_{11}^\top + \bPhi_{11}\bL\bQ_c\bL^\top{\bPhi^{(1)}}^\top\bigr)\tau^4 \nonumber\\
    \text{Order 5:}&\quad +\;\tfrac{8}{15}\bigl(\bL\bQ_c\bL^\top\bPhi_{22}^\top + \bPhi_{22}\bL\bQ_c\bL^\top + \bPhi_{11}\bL\bQ_c\bL^\top\bPhi_{11}^\top\bigr)\tau^5 \nonumber\\
    &\quad +\;\tfrac{3}{10}\bigl(\bL\bQ_c\bL^\top\bPhi_{31}^\top + \bPhi_{31}\bL\bQ_c\bL^\top + \bPhi^{(1)}\bL\bQ_c\bL^\top\bPhi_{21}^\top + \cdots\bigr)\tau^5 \nonumber\\
    \text{Order 6:}&\quad +\;\tfrac{7}{30}\bigl(\bL\bQ_c\bL^\top\bPhi_{41}^\top + \bPhi_{41}\bL\bQ_c\bL^\top + \cdots\bigr)\tau^6 \nonumber\\
    &\quad +\;\tfrac{11}{30}\bigl(\bL\bQ_c\bL^\top\bPhi_{32}^\top + \bPhi_{32}\bL\bQ_c\bL^\top + \cdots\bigr)\tau^6.
    \label{eq:Q-pwl-extra}
\end{align}

The total covariance is assembled by propagating each subsegment's contribution through subsequent transition functions:
\begin{equation}
    \bQ = \sum_{i=0}^{P-2} \left(\prod_{j=i+1}^{P-2}\bPhi_{\mathrm{sub}}^{(j)}\right) \bQ_{\mathrm{sub}}^{(i)} \left(\prod_{j=i+1}^{P-2}\bPhi_{\mathrm{sub}}^{(j)}\right)^\top.
    \label{eq:Q-propagation}
\end{equation}

%==========================================================================
\section{Impact on Estimator Components}
\label{sec:impact}
%==========================================================================

\subsection{Prior Term Assembly}
\label{sec:prior-change}

\paragraph{Original (WNOA).}
The residual~\eqref{eq:e-wnoa} and information matrix~\eqref{eq:qinv-wnoa} are hardcoded.
The function \texttt{assemblePriorTerms} takes no control input argument.

\paragraph{Extended.}
The residual becomes:
\begin{equation}
    \boxed{\be = \bgamma_{k+1} - \bPhi(\mathbf{u}, \Delta s)\,\bgamma_k - \bPhi_{\mathrm{int}}(\mathbf{u}, \Delta s),}
    \label{eq:e-general}
\end{equation}
where $\bPhi$ and $\bPhi_{\mathrm{int}}$ are obtained from \texttt{getTransitionFunction} and \texttt{getTransitionFunctionIntegral} for the control input $\mathbf{u}$ associated with the segment.
The information matrix $\bQ\inv$ is obtained from \texttt{getQInv}.

The Jacobian generalises to:
\begin{equation}
    \bF = \begin{pmatrix}
        -\Jac\inv \Ad & -\bPhi^{[1{:}6,\,7{:}12]} & \Jac\inv & \bzero_6 \\
        -\tfrac{1}{2}\curly{\bvarpi_{k+1}}\Jac\inv \Ad & -\bPhi^{[7{:}12,\,7{:}12]} & \tfrac{1}{2}\curly{\bvarpi_{k+1}}\Jac\inv & \Jac\inv
    \end{pmatrix},
    \label{eq:F-general}
\end{equation}
where $\bPhi^{[1{:}6,\,7{:}12]}$ denotes the top-right $6{\times}6$ block of the transition function matrix, and $\bPhi^{[7{:}12,\,7{:}12]}$ denotes the bottom-right block.
For WNOA, these blocks are $\Delta s\,\bI_6$ and $\bI_6$ respectively, recovering~\eqref{eq:F-wnoa}.

\subsection{Prior Cost Evaluation (Line Search)}
The cost-only function \texttt{getPriorCost} mirrors the same generalisation:
it evaluates $\frac{1}{2}\be^\top\bQ\inv\be$ using the general $\bPhi$, $\bPhi_{\mathrm{int}}$, $\bQ\inv$ without assembling Jacobians.

\subsection{GP Interpolation}
\label{sec:interp-change}

\paragraph{Original (WNOA).}
The matrices $\bPhi$, $\bQ$, $\bQ\inv$ in the GP interpolation formulas~\eqref{eq:Q-interp-wnoa}--\eqref{eq:interp-wnoa} are the hardcoded WNOA expressions.

\paragraph{Extended.}
The interpolation now uses:
\begin{itemize}
    \item $\bPhi(\Delta s)$, $\bQ\inv(\Delta s)$ from \texttt{getTransitionFunction} and \texttt{getQInv} for the full segment.
    \item $\bQ(s_q)$ from \texttt{getQPartial} for the partial interval $[s_k, s_q]$.
    \item $\bPhi(s_q, s_{k+1})$ from \texttt{getTransitionFunctionPartial} for the reverse partial interval.
    \item $\bPhi(s_k, s_q)$ from \texttt{getTransitionFunctionPartial} for the forward partial interval.
    \item $\bPhi_{\mathrm{int}}(\Delta s)$ from \texttt{getTransitionFunctionIntegral} for the full segment (used in input offset computation).
    \item $\bPhi_{\mathrm{int}}(s_q)$ from \texttt{getTransitionFunctionIntegralPartial} for the partial integral.
\end{itemize}

The GP interpolation coefficients become:
\begin{align}
    \bm{\Psi} &= \bQ(s_q)\,\bPhi(s_q, s_{k+1})^\top\,\bQ\inv(\Delta s), \\
    \bm{\Lambda} &= \bPhi(s_k, s_q) - \bm{\Psi}\,\bPhi(\Delta s),
\end{align}
and the interpolated mean accounts for the control input integral:
\begin{equation}
    \bgamma_q = \bm{\Lambda}\,\bgamma_k + \bm{\Psi}\,\bgamma_{k+1} + \bPhi_{\mathrm{int}}(s_q) - \bm{\Psi}\,\bPhi_{\mathrm{int}}(\Delta s).
\end{equation}

\subsection{Initial Guess: PriorMean Type}
\label{sec:priormean}

A new initial guess type \texttt{PriorMean} is added.
Rather than seeding the optimiser with straight rods, it forward-propagates the GP prior mean from the base frame of each robot using the transition functions and control inputs.

Starting from node~0 with $\T_{b_0 i} = \T_{i0}\inv$ and $\bvarpi_0 = \bzero$, each subsequent node is obtained by:
\begin{equation}
    \bgamma_{k+1} = \bPhi(\mathbf{u}_k, \Delta s)\,\bgamma_k + \bPhi_{\mathrm{int}}(\mathbf{u}_k, \Delta s),
\end{equation}
and the global state is recovered via:
\begin{equation}
    \T_{b_{k+1} i} = \exp\!\bigl(\bgamma_{k+1}^{[1:6]\,\wedge}\bigr)\,\T_{b_k i}, \qquad
    \bvarpi_{k+1} = \Jac\!\bigl(\bgamma_{k+1}^{[1:6]}\bigr)\,\bgamma_{k+1}^{[7:12]}.
\end{equation}

This provides a physically motivated initial guess when control inputs are available, substantially improving convergence when the robot shape is far from straight.

%==========================================================================
\section{Summary of Code Changes}
\label{sec:code-changes}
%==========================================================================

\Cref{tab:changes} summarises the structural changes to the header and implementation files.

\begin{table}[h]
\centering
\caption{Summary of code-level changes.}
\label{tab:changes}
\small
\begin{tabular}{@{}p{4.5cm}p{5.5cm}p{5.5cm}@{}}
\toprule
\textbf{Component} & \textbf{Original (WNOA)} & \textbf{Extended (General)} \\
\midrule
\texttt{ControlInput} struct & Not present & \texttt{enum \{None, Const, PwL\}}; values, robot/segment indices \\
\midrule
\texttt{InitialGuessType} enum & \texttt{\{Straight, Last, Custom\}} & \texttt{\{Straight, Last, PriorMean, Custom\}} \\
\midrule
\texttt{constructInitialGuess} & No inputs arg; only straight/custom/last & Accepts \texttt{inputs}; adds \texttt{PriorMean} forward-propagation \\
\midrule
\texttt{assemblePriorTerms} & Hardcoded $\bPhi$, $\bQ\inv$ (WNOA) & Calls \texttt{getTransitionFunction}, \texttt{getQInv}; adds $\bPhi_{\mathrm{int}}$ \\
\midrule
\texttt{getPriorCost} & Hardcoded WNOA & Uses general $\bPhi$, $\bPhi_{\mathrm{int}}$, $\bQ\inv$ \\
\midrule
\texttt{interpolateStates} & Hardcoded WNOA $\bPhi$, $\bQ$, $\bQ\inv$ & Calls partial transition functions \\
\midrule
\texttt{computeStateEstimate} & No inputs forwarding & Passes \texttt{inputs} to initial guess, prior, measurement, interpolation \\
\midrule
\texttt{getTransitionFunction} & Not present & Full/Constant/PwL Magnus expansion \\
\texttt{getTransitionFunctionIntegral} & Not present & Input integral computation \\
\texttt{getQInv} & Not present & Input-dependent covariance \\
\texttt{*Partial} variants & Not present & For GP interpolation sub-intervals \\
\midrule
\texttt{PrecomputedValues} & Not present & Caches $\bPhi$, $\bPhi_{\mathrm{int}}$, $\bQ\inv$ per segment \\
\texttt{m\_LQLT} & Not present & Precomputed $\bL\bQ_c\bL^\top$ \\
\bottomrule
\end{tabular}
\end{table}

%==========================================================================
\section{Discussion}
\label{sec:discussion}
%==========================================================================

\paragraph{Why replace the WNOA prior?}
The WNOA prior assumes that the continuum robot evolves according to a constant-velocity (in strain) stochastic model.
While this is adequate when no prior knowledge of the strain field is available, it becomes suboptimal when:
\begin{itemize}
    \item External loads or actuation produce known strain contributions that can be incorporated as control inputs.
    \item The robot backbone follows a prescribed trajectory in configuration space, making the mean of the GP prior non-trivial.
    \item The physics of the problem (e.g., gravity, tendon tension, magnetic actuation) induces spatially varying curvature that a zero-mean WNOA prior penalises.
\end{itemize}

\paragraph{Backward compatibility.}
When all control inputs are \texttt{None} (the default), the general transition functions reduce to the WNOA expressions.
Specifically:
\begin{itemize}
    \item $\bPhi_{\texttt{None}} = \bPhi_{\mathrm{WNOA}}$,
    \item $\bPhi_{\mathrm{int,None}} = \bzero$,
    \item $\bQ_{\texttt{None}}\inv = \bQ_{\mathrm{WNOA}}\inv$.
\end{itemize}
All existing estimator behaviour is preserved when no control inputs are provided.

\paragraph{Computational cost.}
The Magnus expansion and numerical covariance integration introduce additional computation per segment.
However, the \texttt{PrecomputedValues} cache ensures that each segment's transition function, integral, and covariance are computed only once per optimisation run and reused across prior assembly, cost evaluation, and interpolation.

\paragraph{Numerical considerations.}
The 5th-order truncation of the Magnus expansion is sufficient for the typical arclength intervals encountered in continuum robot discretisation.
For very large intervals or highly curved inputs, the subsegment decomposition in the piecewise-linear mode naturally limits the expansion parameter, maintaining accuracy.

\begin{thebibliography}{9}
\bibitem{lilge2022multi}
T.~Lilge, T.~D.~Barfoot, and J.~Burgner-Kahrs,
``Continuum robot state estimation using Gaussian process regression on $SE(3)$,''
\textit{The International Journal of Robotics Research}, vol.~41, no.~13-14, pp.~1099--1120, 2022.
arXiv: 2210.14842.

\bibitem{lilge2024pcr}
T.~Lilge, J.~Qi, J.~M.~Jackson, T.~D.~Barfoot, and J.~Burgner-Kahrs,
``Parallel continuum robot state estimation using Gaussian process regression on $SE(3)$,''
\textit{IEEE Robotics and Automation Letters}, 2024.
arXiv: 2408.01333.
\end{thebibliography}

\end{document}